\begin{table*}[!t]
\centering
\footnotesize
\renewcommand{\arraystretch}{0.82}
\setlength{\tabcolsep}{3.0pt}
\caption{匹配样本窗口受控稳健性消融实验（$N_{\text{train}}=500$，约对应 21 天遥测数据）：解耦领域先验几何结构与辅助历史数据预算（分位数校准 $r=2, q=1/3$，所有先验组别 $\lambda=1.0$，5 随机种子）。}
\label{tab:information_budget}
\begin{threeparttable}
\begin{tabularx}{\textwidth}{p{2.3cm}p{4.2cm}>{\centering\arraybackslash}p{2.0cm}>{\centering\arraybackslash}p{2.0cm}>{\centering\arraybackslash}p{1.6cm}>{\raggedright\arraybackslash}X}
\toprule
\textbf{风电场} & \textbf{受控实验} & \textbf{模型训练} & \textbf{先验信息} & \textbf{调度成本 ($r=2$)} & \textbf{实证结论} \\
\textbf{与地形特征} & \textbf{组别条件} & \textbf{样本预算} & \textbf{数据预算} & \textbf{(均值 $\pm$ 标准差)} & \textbf{与因果识别分析} \\
\midrule
\multirow{5}{=}{\textbf{Kelmarsh}\\(12.3 MW)\\\textit{开阔平原地形}} 
& A. 无物理先验基准 ($\lambda=0$) & 500 样本 & 无 & $0.09648 \pm 0.00314$ & 数据受限下的纯数据驱动无约束基准 \\
& B. 严格预算匹配真实先验 ($\lambda=1.0$) & 500 样本 & 500 样本 & $0.09499 \pm 0.00307$ & \textbf{领域结构效应} (均值改善 $-1.54\%$)：保留全历史均值增益约 $77\%$ \\
& C. 全历史真实先验 ($\lambda=1.0$) & 500 样本 & 7 个月 (全量) & \textbf{$0.09455 \pm 0.00285$} & 全历史均值增益 ($-2.00\%$)：历史库提供边际增益 ($-0.46\%$) \\
& D. 严格预算匹配打乱先验 ($\lambda=1.0$) & 500 样本 & 500 样本 & $0.13827 \pm 0.02333$ & \textbf{先验特异性}：置换映射导致成本激增 $+43.3\%$ \\
& E. 全历史打乱伪先验 ($\lambda=1.0$) & 500 样本 & 7 个月 (全量) & $0.13798 \pm 0.02338$ & 全历史打乱依然造成严重损害 ($+43.0\%$) \\
\midrule
\multirow{5}{=}{\textbf{Hill of Towie}\\(48.3 MW)\\\textit{复杂高地地形}} 
& A. 无物理先验基准 ($\lambda=0$) & 500 样本 & 无 & $0.10444 \pm 0.00312$ & 样本饥饿下的无约束基准，易发生过拟合 \\
& B. 严格预算匹配真实先验 ($\lambda=1.0$) & 500 样本 & 500 样本 & $0.10095 \pm 0.00305$ & \textbf{领域结构效应} (均值改善 $-3.34\%$)：保留全历史均值增益约 $92\%$ \\
& C. 全历史真实先验 ($\lambda=1.0$) & 500 样本 & 10 个月 (全量) & \textbf{$0.10066 \pm 0.00307$} & 全历史均值增益 ($-3.62\%$)：历史库提供边际增益 ($-0.28\%$) \\
& D. 严格预算匹配打乱先验 ($\lambda=1.0$) & 500 样本 & 500 样本 & $0.15248 \pm 0.00253$ & \textbf{先验特异性}：置换映射导致成本激增 $+46.0\%$ \\
& E. 全历史打乱伪先验 ($\lambda=1.0$) & 500 样本 & 10 个月 (全量) & $0.17209 \pm 0.01703$ & 强约束下先验误设惩罚激增至 $+64.8\%$ \\
\bottomrule
\end{tabularx}
\begin{tablenotes}[flushleft]
\scriptsize
\item \textbf{注}：(1) 所有模型共享相同的 DLinear 架构、优化超参数和后验分位数校准偏移量（$r=2, q=1/3$），并在 2017 年全量测试集和 5 个随机种子（42, 123, 456, 789, 1024）下评估。所有真实与打乱先验组别（B--E组）均在惩罚权重 $\lambda=1.0$ 下执行。(2) A、B、D 组严格共享完全相同的 500 样本训练时间窗口（约 21 天历史）。预算匹配先验保留了全历史先验下观测到的大部分均值改善（两风场分别约保留 77\% 与 92\%）。在 5 个相同种子下的成对差异（Paired Difference）计算中，预算匹配先验相对于 A 组取得显著成对下行风险降低，Kelmarsh 均值成对差异为 $-0.00149$（95\% CI: $[-0.00174, -0.00124]$），Hill of Towie 为 $-0.00349$（95\% CI: $[-0.00368, -0.00331]$）；而将先验扩展至多月全历史数据仅分别贡献了 $-0.00044$（95\% CI: $[-0.00075, -0.00013]$）与 $-0.00029$（95\% CI: $[-0.00041, -0.00016]$）的边际增益。(3) 该成对对比表明，在时间数据窗口匹配下，领域结构化的经验先验效应构成了观测到的大部分风险降低来源，而非单纯源自额外历史数据量的扩充。
\end{tablenotes}
\end{threeparttable}
\end{table*}
